\documentclass[11pt]{article}
\usepackage[margin=1in]{geometry}
\usepackage[T1]{fontenc}
\usepackage{lmodern,amsmath,amssymb,amsthm,mathtools,hyperref}
\setlength{\emergencystretch}{3em}
\newtheorem{theorem}{Theorem}[section]
\newtheorem{lemma}[theorem]{Lemma}
\newtheorem{corollary}[theorem]{Corollary}
\newtheorem{proposition}[theorem]{Proposition}
\newcommand{\E}{\mathbb E}
\newcommand{\tr}{\operatorname{tr}}
\title{Compatible sign words and\\an exact self-dual code pivot}
\author{Joshua Mark\thanks{ORCID: \href{https://orcid.org/0009-0003-5905-2382}{0009-0003-5905-2382}. AI tools assisted with derivation, independent review, computation and exposition.}}
\date{20 November 2026 --- research draft; priority not established}
\begin{document}\maketitle
\section{Compatible codewords and a change of self-dual code}
\label{sec:compatible-codes}

Put $N=4q$ with $q\ge3$ odd. A sign class means a pair $\{v,-v\}$;
use the representative whose first coordinate is positive. Set
$\pi(v)=\prod_jv_j$, which is well-defined on classes because $N$ is even.
We call a family modularly orthogonal if all its pairwise inner products
are divisible by $q$. The following proof requires no primality.

\begin{theorem}[Two-parity rigidity]\label{code:parity}
In a modularly orthogonal family of distinct sign classes, the
classes of each fixed parity are pairwise orthogonal over the reals.
If both parities occur, each parity has at most four classes.
Consequently any such family with at least nine classes is a real
partial Hadamard matrix, and every family has at most $N$ classes.
The mixed-parity bound of eight is sharp for every odd $q\ge3$.
\end{theorem}
\begin{proof}
For distinct nonantipodal signs $u,v$, their inner product is even,
divisible by odd $q$, and strictly between $-4q$ and $4q$. It is
therefore $0$ or $\pm2q$. Since
\[
 \pi(u)\pi(v)=(-1)^{(N-u\cdot v)/2},
\]
equal parities force zero and opposite parities force $\pm2q$.
Within one parity the normalized rows are orthonormal. Against
any normalized row of the opposite parity, every squared inner
product is $1/4$. Bessel's inequality limits the first class to
four rows, and then does the same for the second.
If only one parity occurs there are at most $N$ orthogonal rows;
otherwise there are at most eight, which is less than $N$.

For sharpness start with the four rows of $H_4$. Adjoin the four
antipodal classes of sign vectors of the form
$\frac12\sum_{i=1}^4\sigma_i h_i$ that have all entries $\pm1$.
The two sets are orthogonal bases of $\mathbb R^4$ whose cross
inner products are $\pm2$. Repeat each coordinate $q$ times.
This gives four classes of each parity, with cross products
$\pm2q$ and within-class products zero.
\end{proof}

This is an elementary rigidity result within the classical subject
of modular Hadamard matrices, not a priority claim; see Marrero and
Butson, \emph{Modular Hadamard matrices and related designs} (1973),
\url{https://doi.org/10.1016/0097-3165(73)90073-3}.
The modulus here grows with the target order. Existence theorems
for a fixed small modulus cannot simply supply this hypothesis.

\begin{corollary}[Five rows remove compatibility tests]\label{code:five}
Let $C$ be a submodule of $(\mathbb Z/q\mathbb Z)^N$ with
$C\subseteq C^\perp$, and let
\[
 S(C)=\{v\in\{\pm1\}^N:v\bmod q\in C,\ v_1=1\}.
\]
If $C$ contains five mutually orthogonal sign rows, every pair of
distinct members of $S(C)$ is orthogonal over the reals.
Thus $5\le |S(C)|\le N$.
\end{corollary}
\begin{proof}
The five prescribed orthogonal classes all have the same parity.
One opposite-parity class in $C$ would violate the four-row Bessel
bound. All sign classes in $C$ consequently have this same parity,
and Theorem~\ref{code:parity} applies.
\end{proof}

This changes the task from finding a clique among modular candidates
to counting sign words in one self-orthogonal code. It does not prove
that the count is large. In particular a dual kernel $P^\perp$ of a
short frame is usually not self-orthogonal; this corollary must not
be applied to that entire candidate set.
Five initial orthogonal rows do exist without assuming $H_N$:
write $N=12+8k$ and concatenate the first five rows of a fixed
$H_{12}$ with $k$ copies of the first five rows of $H_8$.
Their modular row span is self-orthogonal, but this construction
alone supplies no useful lower bound beyond five sign classes.

\subsection{A precise restricted-alphabet existence criterion}
The elementary three-column completion theorem proved in
\path{../round5/collision_entropy.tex} says that $N-3$ orthogonal
sign rows of length $N$ can be completed to a full Hadamard matrix.
For $N\ge12$, $N-3\ge9$. The preceding theorem therefore gives:
\[
 \begin{split}
 H_N\text{ exists}\quad\Longleftrightarrow\quad&
 \text{some self-orthogonal submodule }C\subseteq
 (\mathbb Z/q\mathbb Z)^N\\[-2pt]
 &\text{has }|S(C)|\ge N-3.
 \end{split}
\]
For the forward implication take the module generated by the rows
of $H_N$, whose inner products are zero modulo $q$, including the
diagonal $N$. For the reverse implication the sign words are a
partial Hadamard matrix by Theorem~\ref{code:parity}, and complete
its $N-3$ rows. This is an equivalent target, not a count estimate.
One may equivalently require the module to be self-dual. Indeed the
row module of a full $H_N$ is self-dual even for composite odd $q$:
if $Hz=qy$ for integer lifts, then $H^TH=4qI$ implies
$4z=H^Ty$. Since four is invertible modulo $q$, the residue of
$z$ belongs to the row module. Self-orthogonality gives the reverse
inclusion. This is the classical Hadamard-code construction, not
a new self-dual-code existence theorem; see Munemasa and Tamura,
\emph{The codes and the lattices of Hadamard matrices},
\url{https://doi.org/10.1016/j.ejc.2011.11.007}.

Only odd $q$ are needed for the full conjecture: write an arbitrary
multiple of four as $4q\,2^k$ with $q$ odd, and tensor an order-$4q$
matrix with a Sylvester matrix of order $2^k$. The case $q=1$
is already Sylvester. Neither this reduction nor the criterion
establishes existence at the remaining odd multipliers.

\subsection{A self-dual neighbor preserving existing sign rows}
Now assume $q=p$ is an odd prime and work over $\mathbb F_p$.
Let $C=C^\perp$ be self-dual, with at least five orthogonal sign
classes in $S(C)$. If a sign vector $v$ is orthogonal over the
integers to every member of $S(C)$ but does not belong to $C$, put
\[
 C'=(C\cap v^\perp)+\langle v\rangle.
\]
\begin{proposition}[Exact code pivot]\label{code:pivot}
The space $C'$ is self-dual, $\dim(C\cap C')=N/2-1$, and
$S(C)\subsetneq S(C')$. Conversely, if no real sign extension of
$S(C)$ exists, no self-orthogonal code preserving every member
of $S(C)$ can acquire a new sign class.
\end{proposition}
\begin{proof}
The functional $c\mapsto c\cdot v$ on $C$ is nonzero since
$C=C^\perp$ and $v\notin C$, so its kernel has dimension $N/2-1$.
Also $v\cdot v=N=0$ in $\mathbb F_p$. The displayed sum is
self-orthogonal of dimension $N/2$, hence self-dual.
It contains the old sign words because their real inner products
with $v$ vanish, and it contains the new class of $v$.
The assumption $v\notin C$ is automatic for a new real extension:
an extension is not one of the classes being extended, whereas
membership in $C$ would put its normalized representative in $S(C)$.
The converse follows from the same corollary: any acquired class
would be a real extension of all preserved classes.
\end{proof}

Thus repeated extension-guided pivots strictly increase the sign
count and stop after at most $N-|S(C)|$ moves, either at a full matrix
or at a genuinely terminal partial frame. This is not an algorithm
that escapes terminal frames. Moving to a different code can remove
an artificial restriction, but cannot make an incompatible new row
orthogonal to all existing rows.

\subsection{Exact known-order experiments}
The verifier exhausts all self-dual ternary codes containing the
first five rows of a specified Paley $H_{12}$. Their span has rank
five, and there are exactly two self-dual extensions, each with
twelve sign classes. These are results for that prescribed seed,
not a classification of all five-row frames.

At order twenty, a deterministic finite-field construction extends
the first five rows of a Paley $H_{20}$ to a self-dual code over
$\mathbb F_5$. Exhausting its $2^9$ normalized sign assignments in
systematic coordinates yields only eight classes. This disproves
the suggestion that any self-dual extension containing five rows
automatically gives a full matrix. The eight-row frame has rank
eight modulo five. Enumerating its modular kernel gives forty
classes: the eight existing rows and thirty-two exact extensions.
Using the first extension in a single pivot gives
\[
             |S(C)|=8\quad\longrightarrow\quad |S(C')|=20.
\]
All twenty rows are checked pairwise over the integers.
The finite-field bases and the extension witness are recorded in
\path{codes/checks.json}.

Finally self-duality alone need not give even one sign word.
Over $\mathbb F_3$, the length-four code generated by
$(1,0,1,1)$ and $(0,1,1,-1)$ is self-dual and has no word of full
support. Its threefold direct sum is a self-dual length-twelve code
with no sign words. A useful general lower bound must use more than
dimension and self-duality. No count near $N$ has been proved at
unknown orders.
\section{A terminal code trap with an optimal two-row escape}

This certificate concerns order twenty, where Hadamard matrices are
already known. It gives an exact trap for monotone row addition and
an exact minimal-loss escape in the self-dual-code construction.

\begin{theorem}[Finite certified example]
There is a self-dual ten-dimensional code $C\subseteq\mathbb F_5^{20}$
with exactly twelve antipodal classes of sign words. Representatives
with first coordinate $+1$ form a real orthogonal $12\times20$ sign
frame $P$, and span $C$ over $\mathbb F_5$.
No further sign row is orthogonal to every row of $P$.
Every new sign class is orthogonal to at most ten old rows, so no
one-row deletion admits even one genuinely new row.
Deleting the two old rows indexed $3,8$ (zero based) permits a full
Hadamard completion of the remaining ten rows. Its mod-five code
$C'$ is a self-dual neighbor: $\dim(C\cap C')=9$.
One explicit neighbor pivot takes the sign count from twelve to twenty.
Thus two is the exact minimum number of old rows lost in such an escape.
\end{theorem}

The complete basis and twelve sign representatives are stored in
\texttt{search.json}. The construction starts with the first five rows
of the Paley order-twenty matrix and extends their span over
$\mathbb F_5$ using deterministic seed $2026112000$.
The construction is only a way to find the certificate; correctness
is checked from its stored entries. The code basis is self-orthogonal
of rank ten. Enumerating its $2^9=512$ systematic sign assignments
finds exactly twelve normalized sign words, of mod-five rank ten;
their real Gram matrix is $20I_{12}$.
If a new sign row were exactly orthogonal to $P$, its residue would
belong to $C^\perp=C$, hence to the already exhausted twelve classes.
Every such class fails orthogonality to itself, proving terminality.

An independent binary enumeration checks all $2^{19}=524288$
normalized sign words using only integer Hamming distances. Excluding
the twelve existing words, its compatibility counts are
\[
\begin{array}{c|rrrrrrrrrrr}
\text{compatible}&0&1&2&3&4&5&6&7&8&9&10\\\hline
\text{new classes}&262840&6840&27092&48948&67344&58976&34920&12600&3992&616&108
\end{array}
\]
No new word is orthogonal to eleven or twelve old rows. The 108
best words group by their two nonorthogonal old rows. For the pair
$(3,8)$, all ten available words are mutually orthogonal; together
with the ten retained old words they give Gram matrix $20I_{20}$.
Their normalized negative-coordinate bit masks are
\[
7676,126672,307898,444300,581982,645928,683178,816740,855582,886242.
\]
Bit $j$ represents coordinate $j$, with bit one denoting a minus sign.
The retained old rows have modular rank nine, the completed code
rank ten, and the union of the old and new code bases rank eleven.
These prove the claimed intersection dimension. If $v$ is the first
new row, then
\[
 C'=(C\cap v^\perp)+\langle v\rangle,
\]
as independently checked by row reduction.

The scripts \texttt{search.py} and \texttt{verify.py} use only the
Python standard library and original project code from the adjacent
\texttt{codes/check.py}. The latter also writes
\texttt{verification.json} and \texttt{exchange.json}. No external
implementation is imported. This example disproves the strategy of
always retaining all previously selected rows; it does not assert
that two-row neighbor escapes suffice in general or yield unknown
Hadamard orders.

\begin{theorem}[Elementary tensor persistence]
If a sign frame $P$ admits no orthogonal sign row and $H$ is a square
Hadamard matrix of order $t$, then $P\otimes H$ also admits no
orthogonal sign row. In particular this certificate gives terminal
$12t\times20t$ frames whenever such an $H$ is supplied.
\end{theorem}
\begin{proof}
Reshape a hypothetical new sign vector into a $20\times t$ sign
matrix $X$. Orthogonality to $P\otimes H$ says $PXH^T=0$, under
the corresponding vectorization convention. Since $H$ is invertible,
$PX=0$. Each column of $X$ would therefore be a sign vector
orthogonal to $P$, a contradiction.
\end{proof}
\section{An exact local count identity and a two-code decomposition}
\label{sec:neighbor-counts}

The preceding pivot suggests counting over code neighbors, rather
than asking for an extension preserving every row. The following
identity makes that approach quantitative. Here $p$ is an odd prime,
$N=4p$, $m=N/2$, $C=C^\perp$, and $G\subsetneq C$ has dimension
$s=m-h$, $h\ge1$. Put
\[
 r_C=|S(C)|,\quad b_G=|S(G)|,\quad
 M_G=|\{v\in\{\pm1\}^N:v_1=1,\ v\bmod p\in G^\perp\}|.
\]
Let $\mathcal N_G(C)$ be the self-dual codes $C'\ne C$ with
$\dim(C\cap C')=m-1$ and $G\subseteq C'$.

\begin{theorem}[Exact restricted-neighbor incidence count]
\label{code:neighbor-average}
There are $(p^h-1)/(p-1)$ codes in $\mathcal N_G(C)$, and
\[
 \boxed{\quad
 \sum_{C'\in\mathcal N_G(C)}\bigl(|S(C')|-r_C\bigr)
 =M_G-r_C-(r_C-b_G)p^{h-1}.\quad}
\]
In particular, if the right-hand side is positive, some neighbor
strictly increases the sign count while preserving $G$.
\end{theorem}
\begin{proof}
A hyperplane $D$ of $C$ containing $G$ has $D^\perp/D$ a nondegenerate
two-dimensional hyperbolic space. It has exactly two isotropic lines,
one giving $C$ and the other a unique neighbor $C'$.
To see the last assertion directly, choose isotropic $e$ on the first
line and $f$ with $e\cdot f=1$, then replace $f$ by
$f-(f\cdot f)e/2$. In the resulting basis the norm is $2xy$,
so the isotropic lines are precisely the two coordinate axes.
Counting hyperplanes of $C/G$ gives the stated neighbor count.

Each normalized sign vector outside $C$ but in $G^\perp$ is isotropic,
since its norm is $N=0$ modulo $p$. It belongs to exactly one of
these neighbors, namely
$(C\cap v^\perp)+\langle v\rangle$. Thus the total number of
new sign classes, counted over the neighbors, is $M_G-r_C$.
An old sign class in $G$ is never lost. An old class outside $G$
is excluded by precisely $p^{h-1}$ hyperplanes containing $G$,
the difference between the hyperplane counts in dimensions
$h$ and $h-1$. Summing gains minus losses proves the identity.
\end{proof}

This is a finite incidence identity in the classical geometry of
self-dual code neighbors. It is not a general positivity theorem.
In particular, at $G=0$ its right-hand side is
$2^{N-1}-r_C(1+p^{m-1})$, which can be negative even when the
code has far fewer than $N$ sign classes. Unrestricted averaging
therefore does not guarantee a useful move.

\begin{corollary}[A modular pool splits into two codes]
\label{code:two-codes}
If $G$ has dimension $m-1$, there is exactly one other self-dual
code $C'$ containing $G$, and
\[
 \{ \text{sign classes in }G^\perp\}=S(C)\cup S(C'),\qquad
 S(C)\cap S(C')=S(G),\qquad
 |S(C')|=M_G-r_C+b_G.
\]
If $G$ contains five orthogonal sign rows, the sign words of each
code separately form a real orthogonal frame.
\end{corollary}
\begin{proof}
An isotropic residue in $G^\perp$ projects to one of the two
isotropic lines of $G^\perp/G$, hence belongs to $C$ or $C'$.
Their intersection is $G$. Count the union, and apply
Corollary~\ref{code:five} for real orthogonality.
\end{proof}

This removes a compatibility search in the codimension-one case.
For the terminal twelve-row order-twenty certificate, the ten rows
retained after deleting indices $3,8$ span such a $G$, and exact
enumeration gives
\[
 r_C=12,\qquad b_G=10,\qquad M_G=22,\qquad |S(C')|=20.
\]
The twelve exact extensions of the retained frame consist of the two
deleted old classes and the ten new classes. The two-code theorem,
rather than an independent guess that all candidates are compatible,
explains the resulting completion.

The checker also verifies the incidence identity for all four
neighbors containing a specified four-dimensional subspace of a
self-dual ternary code at order twelve. These are finite examples.
An all-order lower bound on the modular pool, or a guarantee that
some choice of retained space makes the displayed surplus positive,
has not been established.

\end{document}
